\section*{Chapter \ref{ch:m2:thinking-in-expected-value} --- Thinking in Expected Value}

\begin{solution}{exo:m2:thinking-in-expected-value:1}
Edge $0.4 \times 2 - 0.6 = 0.2$ per unit staked; Kelly fraction $0.2/2 = 0.1$.
\end{solution}

\begin{solution}{exo:m2:thinking-in-expected-value:2}
It is dominated by a few enormously lucky paths: the expected wealth grows like $1.04^{300}$
while almost every player's wealth grows like $e^{300g}$, which is far smaller. The median,
or the \omterm{def:m2:thinking-in-expected-value:growth}{growth rate}, describes what a typical player gets.
\end{solution}

\begin{solution}{exo:m2:thinking-in-expected-value:3}
$2c - c^2$: 75\% at half Kelly, 43.75\% at a quarter.
\end{solution}

\begin{solution}{exo:m2:thinking-in-expected-value:4}
$0.2^{2/1 - 1} = 20\%$ at full Kelly; $0.2^{3} = 0.8\%$ at half Kelly.
\end{solution}

\begin{solution}{exo:m2:thinking-in-expected-value:5}
$c = 2/(1 + \ln 0.1/\ln 0.8) = 0.18$: less than a fifth of Kelly.
\end{solution}

\begin{solution}{exo:m2:thinking-in-expected-value:6}
$t^2/(t^2 + 1) = 4/5 = 0.8$.
\end{solution}

\begin{solution}{exo:m2:thinking-in-expected-value:7}
About 94\% at half Kelly and at Kelly (94.2\% and 93.7\% in 4\,000 seeded players), close to
the 95\% the authors quote for constant fractions of 10 to 20\%.
\end{solution}

\begin{solution}{exo:m2:thinking-in-expected-value:8}
With a Sharpe ratio of 0.1 over 200 trades the edge has a $t$-statistic of about 1.4: it is
uncertain, and staking the full estimated Kelly fraction overbets whenever the estimate is
high. The expected growth is maximised at about two thirds of the estimated fraction,
and full Kelly also carries a 50\% chance of ever halving the capital.
\end{solution}

\begin{solution}{pb:m2:thinking-in-expected-value:1}
\textbf{1.} Edge 0.1 per unit staked, Kelly fraction 0.1, growth 0.501\% per trade.
\textbf{2.} Mean 0.1, standard deviation 0.995, Sharpe ratio 0.1005 per trade.
\textbf{3.} 50\%.
\textbf{4.} 0.46 of Kelly.
\textbf{5.} To second order the \omterm{def:m2:thinking-in-expected-value:growth}{growth rate} at twice Kelly is zero: the gain from the edge
is exactly offset by the variance drag of the larger stakes.
\textbf{6.} The mean is estimated with a standard error of 0.0704; $t = 1.42$.
\textbf{7.} $c^* = 0.669$.
\textbf{8.} 0.7, the best of the multiples tried, in steps of 0.1.
\textbf{9.} 19\%: 29.1 against 35.9 basis points a trade.
\textbf{10.} $0.5^{2/0.669 - 1} = 25\%$.
\textbf{11.} With 50 trades, 0.34; with 2\,000, 0.95.
\textbf{12.} Growth falls quadratically around the Kelly fraction but the drag grows with
the square of the stake, and an overestimated edge raises the stake; the losses from
overbetting include a real chance of ruin, those from underbetting only slower growth.
\textbf{13.} Lower it: losses larger than the variance suggests make large stakes costlier
than the quadratic approximation shows.
\textbf{14.} Take the smaller of the shrunk Kelly multiple and the multiple that keeps the
\omterm{def:m2:thinking-in-expected-value:frac}{risk of ruin} within the tolerance.
\textbf{15.} That people with a clear edge, when free to size, overbet and bet erratically:
a sizing rule decided in advance protects the edge from the trader.
\textbf{16.} Because edges are estimated with error and decay, returns have fat tails,
strategies are correlated, and investors and counterparties impose \omterm{def:m2:thinking-in-expected-value:frac}{drawdown} limits long
before ruin.
\textbf{17.} Bayesian updating of the edge's estimate and its uncertainty after each
trade, recomputing the shrunk fraction, with limits on how fast it may change.
\textbf{18.} When the goal is not long-run growth of one bankroll: a fixed horizon, a cap
on payout, a hard \omterm{def:m2:thinking-in-expected-value:frac}{drawdown} limit, or utility with more risk aversion than the logarithm.
\textbf{19.} Named result: \emph{the uncertain edge}: with the edge estimated from 200 trades
at a Sharpe ratio of 0.1, the best stake is two thirds of the estimated Kelly fraction
(0.669 by the formula, 0.7 in simulation), which earns 19\% more growth than full Kelly.
\textbf{20.} Because the edge is an estimate, and overbetting it costs more than
underbetting.
\end{solution}

\begin{solution}{iq:m2:thinking-in-expected-value:1}
Twenty per cent: the Kelly fraction $p - q/b = 0.6 - 0.4$, which maximises the long-run
\omterm{def:m2:thinking-in-expected-value:growth}{growth rate} at about 2\% a flip. Betting more raises the variance drag faster than the
gain; above about 39\% wealth shrinks in the long run.
\iqlookfor{edge over \omterm{def:m2:thinking-in-expected-value:edge}{odds}, and why not more.}
\end{solution}

\begin{solution}{iq:m2:thinking-in-expected-value:2}
Expected wealth rewards staking everything, since a few huge outcomes dominate it;
expected log wealth is the \omterm{def:m2:thinking-in-expected-value:growth}{growth rate} that almost every path achieves over many bets,
and it penalises variance. Maximising it gives Kelly; maximising the mean gives ruin.
\iqlookfor{typical versus average outcome.}
\end{solution}

\begin{solution}{iq:m2:thinking-in-expected-value:3}
Over a short period $dt$, the log return of wealth with fraction $f$ in the asset (and
the rest at zero rate) has mean $f\mu\,dt - \tfrac12 f^2\sigma^2 dt$; maximising gives $f^* =
\mu/\sigma^2$ and growth $\mu^2/(2\sigma^2)$, half the squared Sharpe ratio.
\iqlookfor{Ito's correction and the quadratic.}
\end{solution}

\begin{solution}{iq:m2:thinking-in-expected-value:4}
Half Kelly keeps three quarters of the growth with half the volatility; it guards against
overestimated edges and fat tails; and it keeps \omterm{def:m2:thinking-in-expected-value:frac}{drawdowns} within what investors and
risk limits tolerate: the chance of ever halving falls from 50\% to 12.5\%.
\iqlookfor{growth versus risk, and estimation error.}
\end{solution}

\begin{solution}{iq:m2:thinking-in-expected-value:5}
The Kelly weights are $\mu_i/\sigma_i^2 = \mathrm{SR}_i/\sigma_i$: with equal volatilities,
twice the capital to the Sharpe-2 strategy; in general in proportion to Sharpe ratio over
volatility, then scaled down for estimation error and limits.
\iqlookfor{$\Sigma^{-1}\mu$ with no correlation.}
\end{solution}

\begin{solution}{iq:m2:thinking-in-expected-value:6}
Each day, estimate each strategy's edge and risk from its live and backtest history with
their uncertainty; compute the shrunk Kelly fraction, cap it by the \omterm{def:m2:thinking-in-expected-value:frac}{drawdown} rule, by
correlations with the other strategies and by capacity; smooth changes; publish the
allocations with their inputs; override on breaches; and log everything for review.
\iqlookfor{estimation, shrinkage, limits, smoothing, audit.}
\end{solution}
